%%%%%%%%%%%%%%%%%%%%%%%%%%%%%%%%%%%%%%%%%%%%%%%%%%%%%%%%%%%%%%%%%%%%%%%%%%%%%%%%%%%%%
%
%  PLANTILLA DE TRABAJO DE GRADO
%  Universidad de San Buenaventura
%  By GHS & PADIA - 2026-2
%
%%%%%%%%%%%%%%%%%%%%%%%%%%%%%%%%%%%%%%%%%%%%%%%%%%%%%%%%%%%%%%%%%%%%%%%%%%%%%%%%%%%%%
%%%%%%%%%%%%%%%%%%%%%%%%%%%%%%%%%%%%%%%%%%%%%%%%%%%%%%%%%%%%%%%%%%%%%%
%  Fuente del documento
%
%  Por defecto se usa TeX Gyre Termes: un clon libre de Times New Roman con
%  las mismas metricas (ocupa exactamente el mismo espacio en la pagina) que
%  viene incluido en todas las distribuciones de LaTeX y en Overleaf.
%
%  Se carga por nombre de archivo, no por nombre de fuente del sistema. Eso
%  es intencional: buscar fuentes del sistema obliga a XeLaTeX a recorrer el
%  catalogo de fuentes de la maquina, lo que en Overleaf puede tardar tanto
%  que la compilacion supera el limite de tiempo del plan gratuito.
%
%  Si trabaja en Windows o macOS y prefiere la Times New Roman autentica,
%  ponga el interruptor de abajo en \timesnewromantrue. En Overleaf dejelo
%  como esta.
%
%  La plantilla funciona con XeLaTeX, LuaLaTeX y pdfLaTeX: la fuente se
%  ajusta sola segun el motor.
%%%%%%%%%%%%%%%%%%%%%%%%%%%%%%%%%%%%%%%%%%%%%%%%%%%%%%%%%%%%%%%%%%%%%%

%  <<<<  interruptor: cambie a \timesnewromantrue para usar Times New Roman
\newif\iftimesnewroman
\timesnewromanfalse

\usepackage{iftex}

% verdadero con XeLaTeX o LuaLaTeX; falso con pdfLaTeX
\newif\ifmotormoderno
\ifXeTeX  \motormodernotrue \fi
\ifLuaTeX \motormodernotrue \fi

\ifmotormoderno
    %------------------------------------------------------------------
    %  XeLaTeX o LuaLaTeX
    %------------------------------------------------------------------
    \usepackage[no-math]{fontspec} % quite [no-math] si presenta errores

    \iftimesnewroman
        % fuentes instaladas en el sistema (Windows y macOS)
        \setmainfont{Times New Roman}
        \setmonofont{Courier New}[Scale=MatchLowercase]
    \else
        % clones libres, cargados por nombre de archivo (rapido, sin
        % consultar el catalogo de fuentes del sistema)
        \setmainfont{texgyretermes}[
            Extension      = .otf,
            UprightFont    = *-regular,
            BoldFont       = *-bold,
            ItalicFont     = *-italic,
            BoldItalicFont = *-bolditalic,
        ]
        \setmonofont{texgyrecursor}[
            Extension      = .otf,
            UprightFont    = *-regular,
            BoldFont       = *-bold,
            ItalicFont     = *-italic,
            BoldItalicFont = *-bolditalic,
            Scale          = MatchLowercase,
        ]
    \fi
\else
    %------------------------------------------------------------------
    %  pdfLaTeX: no puede usar fuentes del sistema, siempre usa los clones
    %------------------------------------------------------------------
    \usepackage[T1]{fontenc}
    \usepackage[utf8]{inputenc}
    \usepackage{tgtermes}          % clon libre de Times New Roman
    \usepackage{tgcursor}          % clon de Courier, para el codigo
    \renewcommand{\familydefault}{\rmdefault}
\fi
